\section{Implementation and Module-Wise Coding}
\label{ch:impl}

\subsection{Development Environment}
The system is written in Python~3.10 with \texttt{asyncio}, the \texttt{cryptography} library and SQLite; the API uses FastAPI, uvicorn and pydantic; the web application uses React~18, Vite and Tailwind; testing uses pytest and pytest-asyncio; version control is Git on GitHub. The source is about 6\,500 lines of Python (\texttt{mesh/}, \texttt{api/}, \texttt{sim/}, \texttt{story/}), 4\,500 lines of tests and 3\,100 lines of React.

\subsection{Overall Implementation Strategy}
The build order followed what blocks what. \textbf{Phase~1:} the 52-byte packet and the \texttt{Transport} interface, so every track could work against a stable boundary. \textbf{Phase~2:} discovery, single-hop relay and the simulated transport, so that routing and reliability could be tested in one process. \textbf{Phase~3:} routing, reliability, security and storage, wired into \texttt{Node}. \textbf{Phase~4:} the API and React application, the scenario runner and metrics, engine fixes found by integration, and the multi-device deployment.

\subsection{Module-Wise Implementation}
\label{sec:modules}
Table~\ref{tab:impl} gives, for each module, its code reference (file and size), its purpose, input and output, processing logic, dependencies and test evidence.
\begin{table}[!htbp]
\caption{Modules: code reference, logic, dependencies and tests}
\label{tab:impl}
\footnotesize
\begin{tabularx}{\textwidth}{@{}>{\raggedright\arraybackslash}p{3.55cm}>{\raggedright\arraybackslash}X>{\raggedright\arraybackslash}p{2.25cm}>{\raggedright\arraybackslash}p{1.8cm}@{}}
\toprule
\textbf{Module (lines)} & \textbf{Purpose, input $\rightarrow$ output, processing, details} & \textbf{Depends on} & \textbf{Tests}\\ \midrule
\texttt{mesh/core/packet.py} (362) & wire format, \texttt{Packet} $\leftrightarrow$ bytes. \texttt{struct} format asserted to be 52 bytes; every field range-checked; \texttt{unpack} rejects bad magic and truncation of each region; \texttt{remaining\_length} frames stream reads; \texttt{identity\_bytes} is the signed content & -- & 47 + 4\,000 fuzz\\
transports (461) & bytes between neighbours: simulated (seeded loss, latency), multicast/TCP, and the multi-device \texttt{LanTransport} (Section~\ref{sec:lan}) & \texttt{packet} & 15 + 6\\
\texttt{discovery.py} (292) & HELLOs $\rightarrow$ neighbour table, peer events, link quality over a 20-slot window reported only when it moves by 0.1; \texttt{visible\_neighbours} range emulation & \texttt{transport} & 15 + 12\\
\texttt{relay.py} (153) & received packets $\rightarrow$ deliver, forward or drop (algorithm A4); one drain task over a priority queue & \texttt{reliability} & 18 + 3\\
\texttt{routing.py} (546) & reports $\rightarrow$ graph, next hop. \texttt{MeshGraph}, \texttt{edge\_cost}, sequence-numbered gossip with quality, failures, role and key, cached Dijkstra, \texttt{RESCUE\_\_} terminals, hop-count switch for the comparison & \texttt{discovery} & 14 + 14 + 33\\
\texttt{reliability.py} (240) & \texttt{DedupSet} (300\,s), \texttt{MessagePriorityQueue}, \texttt{AckTracker} (3 attempts), \texttt{StoreForwardQueue} (500 messages, 1800\,s) & \texttt{packet} & 18 + 20\\
\texttt{crypto.py} (131) & AES-256-GCM with nonce prefix and associated data; Ed25519 sign and non-raising verify; one exception type & \texttt{cryptography} & 19\\
\texttt{store.py} (258) & tables, WAL, upserts, event queries, \texttt{events\_since}, listing and counting for metrics & \texttt{sqlite3} & 23\\
\texttt{node.py} (513) & composition root: \texttt{send()} encrypts, signs, logs, dispatches; receive verifies, decrypts, logs, ACKs; ACK loop records the failure, recomputes, resends; flush loop expires and retries & all above & 27\\
\texttt{sim/}, \texttt{api/lab.py}, \texttt{api/metrics.py} (3\,000) & simulator with a virtual clock and seven scenarios; metrics read from node logs & \texttt{node} & 5 + 12\\
\bottomrule
\end{tabularx}
\end{table}

\subsection{Front-End Implementation}
The front end is one app: the \textbf{device screen} (\texttt{story/node.html}), a single page served by every device that becomes a \emph{survivor phone}, a \emph{gateway} view or a \emph{rescue console} according to the device's role. A phone is only a screen: the engine runs on the laptop and the phone opens the page over the same Wi-Fi or hotspot. The survivor screen is built to be read under stress: a plain status line, need buttons (need help, trapped, injured, fire or smoke, water or food), a location field, an SOS button held about 1.2\,s to prevent accidental sends, an ``I am safe'' check-in that never delays an SOS, a form to send a text message to a named device, the rescue team's reply shown above a plain-language progress tracker, and the nearby phones with signal strength. The rescue console shows each request with its need, location, route and age; a separate Marked safe list; filters by type and need; a reply action with quick replies; message history with priority and direction filters; CSV export of the received requests; counters; and a network map that shows link quality and highlights the route of a new message. Every screen also has a live panel, ``What this device is doing'', that words the device's own event log in plain English, so discovery and routing can be watched happening. Screens are shown in Section~\ref{sec:shots}. The \textbf{simulation lab} (\texttt{ui/}, a React PWA served by \texttt{api/server.py}) is no longer part of the demonstration. It stays in the repository because it produced the experiments: it runs many real nodes over the simulated radio, switches nodes off and on, edits links, and offers pages for routing, reliability, the wire protocol, security and the simulator.

\subsection{Back-End Implementation}
The back end is FastAPI with a WebSocket. The engine server (\texttt{api/server.py}) exposes the simulation lab (snapshot, reset, kill and revive, link edits, bursts), a node's view, graph and route (weighted and hop-count), \texttt{send} and \texttt{sos}, the rescue view, statistics, one message's full story, the event log, a protocol inspector, a security playground and attack drill, background scenario jobs, and \texttt{/ws}, which streams every logged event. Each device server (\texttt{story/node.py}) exposes \texttt{/api/me} (everything its screen needs), \texttt{/api/sos} (the need and location travel as structured fields), \texttt{/api/send} (a text message to a named device), \texttt{/api/safe} (an ``I am safe'' check-in at the lowest priority), \texttt{/api/events} (the event log in plain English), \texttt{/api/history} (messages filtered by priority and direction), \texttt{/api/export.csv} (the received requests, rescue device only) and \texttt{/api/power}; started with \texttt{--phone} it serves its screen on every network interface so a phone on the same Wi-Fi can open it. Request bodies are validated with pydantic, and engine errors map to clear HTTP errors.

\subsection{Database Implementation}
The three tables of Section~\ref{sec:db} are created with \texttt{CREATE TABLE IF NOT EXISTS}, WAL mode is set at connection time, and the two indexes are created with the schema. Queries used include retained messages by status (\texttt{SELECT * FROM messages WHERE status = ?}), a message's timeline (\texttt{SELECT * FROM events WHERE msg\_id = ? ORDER BY ts\_ms}), event counters (\texttt{SELECT COUNT(*) FROM events WHERE kind = ?}), and the upsert \texttt{INSERT \ldots ON CONFLICT(msg\_id) DO UPDATE}. Every write commits immediately.

\subsection{Communication and Concurrency Implementation}
\label{sec:lan}
The original \texttt{UdpTcpTransport} could not carry a complete \texttt{Node}: it never learned a HELLO sender's address, and it reported the packet's \emph{original} source instead of the previous hop, which breaks link-state flooding. \texttt{lan\_transport.py} fixes both. A JSON \textbf{topology file} lists every device (role, host, port) and the links (who can hear whom); a device talks only to its configured neighbours, which is the range emulation of the honest statement and satisfies DIS-7, so no multicast is needed. Each connection starts with the previous hop's identity followed by the unmodified packet. A \textbf{power switch} makes a device neither send nor hear, like a dead radio. \texttt{python -m story.node --id L1} starts one device (engine plus a small web server); five of them, started together on one machine or one per laptop, form the mesh, each with its own SQLite log, port and screen. In the two-laptop demonstration L1--L3 ran on one laptop and L4--L5 on the other, joined through a phone hotspot (Figure~\ref{fig:deploy}). Recovery is layered: TCP per hop; ACK, retry and re-route end to end; store-and-forward when retries run out.

\subsection{Git Repository Setup}
\label{sec:git}
\textbf{Repository:} \url{\RepoURL}. \textbf{Structure:} \texttt{mesh/core/} (engine), \texttt{api/}, \texttt{sim/}, \texttt{story/}, \texttt{ui/}, \texttt{tests/}, \texttt{docs/}, \texttt{config/}, with README, CONTRIBUTING and requirements from the first phase. \textbf{Branching strategy:} one branch per phase, each inheriting the previous and never merged back (\texttt{develop}, \texttt{develop\_2}, \texttt{develop\_3}, \texttt{develop\_4}, \texttt{develop\_5}). \textbf{Commit conventions:} Conventional Commits (\texttt{feat(routing): \ldots}, \texttt{fix(reliability): \ldots}). \textbf{Contribution of members:} Miti Shah built discovery, the relay, the node skeleton and the simulated transport in Phase~2 and, in Phase~3, Dijkstra pathfinding with gateway and rescue handling and route caching, the dedup set and priority queue, ACK tracking with retry and the store-and-forward queue, AES-256-GCM encryption with Ed25519 signing, and the SQLite log; Twissha Shah built the weighted cost function and link-state gossip, the wiring of routing, reliability, crypto and storage into the node, the multi-device program and device screens, the robustness and socket tests, and the evidence scripts; Tanvi Sheth built the packet format, transports, the API, the React application and the scenario runner, and the Phase~4 engine fixes found by integration. Nine of Miti's Phase~3 commits were made from Twissha's laptop and were first recorded under Twissha's Git account; they were re-attributed to Miti on 24 October 2026 with the agreement of all three members, and the correction is recorded in \texttt{docs/contribution-log.md}. \textbf{Evidence of development history:} Table~\ref{tab:commits}.

\begin{table}[H]
\caption{Commits per team member, read from \texttt{git log --all} (45 commits)}
\label{tab:commits}
\small
\begin{tabular}{@{}lrr@{}}
\toprule
\textbf{Member} & \textbf{Commits} & \textbf{Share}\\ \midrule
Twissha Shah & 24 & 53\,\%\\
Tanvi Sheth & 19 & 42\,\%\\
Miti Shah & 2 & 4\,\%\\
\bottomrule
\end{tabular}
\end{table}
\textbf{Conventional Commits.} 40 of 45 commits (89\,\%) follow the \texttt{type(scope): subject} convention of \texttt{CONTRIBUTING.md}.

\textbf{Branches:} \texttt{main}, \texttt{develop}, \texttt{develop\_2}, \texttt{develop\_3}, \texttt{develop\_4}. \textbf{Tags:} \texttt{v0.3-phase3}.


\subsection{Integration of Modules}
\texttt{Node} registers one dispatcher on the transport that feeds each packet to discovery (HELLO), gossip (LINK\_STATE) and relay (DATA, ACK); each ignores types it does not own. Callbacks connect the pieces: peer-up sets a one-hop route and a graph edge and logs the event; a topology change invalidates cached routes; relay hooks write the message and event log; the ACK loop reports failures to gossip so the cost reacts at once. The API reads the same objects and database, so nothing on screen is computed separately from the engine.

\subsection{Implementation Challenges and Resolutions}
The significant difficulties (Table~\ref{tab:defects} gives each cause and retest) were: a single receive callback that let only one of discovery and relay hear packets; topology known only two hops out, with dead nodes still routable; signatures that failed after the first relay (fixed by signing identity fields only); ACKs and retries lost at relays; a hostile packet raising \texttt{UnicodeDecodeError}; a hang in \texttt{Node.stop()} on Windows; and a transport that could not run a full node (fixed by \texttt{LanTransport}).
